\section{Empirical design and data}\label{sec:empirical-design}

The information-certified portfolio in \Cref{sec:portfolio-choice} is the
theoretical class of allocations that minimize the certified variance bound.
The empirical exercise implements its zero-slack, standardized-asset
specialization; throughout the empirical sections, the resulting canonical
implementation is the news-only allocation.
At this boundary, the carrier scale changes the size of the variance certificate
but not the normalized allocation selected from the observed information
geometry.
With unit marginal scales, $A(x)=1$ and $q=x$, so minimizing the certified
variance objective is equivalent to maximizing the pairwise certificate over
normalized weights.
Implementing the general raw-capital objective would instead require the
marginal scales in \eqref{eq:p3-risk-weights}.
The design next defines the reference populations used to evaluate the
standardized allocation and constructs the news geometry that enters the
decision.
Returns assess the news-only allocation only after the decision has been made.

\subsection{Directional hypothesis and reference-population estimand}
\label{sec:p3-reference-populations}

The directional hypothesis is that the news-only allocation lies in
the lower tail of conventional variance among feasible portfolios with
comparable concentration.
The theory motivates this direction under the maintained return bridge, but it
does not determine a percentile or imply out-of-sample performance.
The news-only allocation is
\[
  q^{\mathrm{news}}\in
  \operatorname*{arg\,max}_{q\ge0,\,\1^\top q=1}
  \frac12q^\top W_2^2q.
\]
It is constructed from the observed W$_2$ geometry before the return covariance
enters the evaluation.

Let
\[
  V(q):=q^\top\widehat\Sigma_R q
\]
be standardized full-sample return variance.
For each reference population $\mathcal G_k$, define the lower-tail
percentile
\[
  p_k:=\Pr_{q\sim\mathcal G_k}
  \{V(q)\le V(q^{\mathrm{news}})\},
  \qquad
  \widehat p_k:=B^{-1}\sum_{b=1}^{B}
  \1\{V(q_k^{(b)})\le V(q^{\mathrm{news}})\},
  \qquad
  B=\ppnum[0]{reference-draws}.
\]
The estimand $p_k$ is the share of portfolios from reference population
$\mathcal G_k$ whose variance is no greater than the news-only allocation's.
Accordingly, a smaller $\widehat p_k$ means that fewer reference portfolios
achieve equally low or lower variance, placing the candidate nearer the bottom
of the conventional-variance distribution.
This quantity is a Monte Carlo estimate of a descriptive reference-population
percentile, not a classical p-value, and the analysis imposes no rejection
threshold.

Each $\mathcal G_k$ starts with a symmetric Dirichlet draw on the long-only
simplex, clips every weight at its cap, and renormalizes until the allocation
is feasible.
The four $(\alpha,\mathrm{cap})$ pairs are
\[
  (1,12.5\%),\qquad (1,15\%),\qquad
  (0.8,15\%),\qquad (0.5,15\%).
\]
The first two populations isolate cap sensitivity.
The effective number of names is
\[
  N(q):=\left(\sum_i q_i^2\right)^{-1}.
\]
The $\alpha=0.8$ law supplies a fixed, moderately concentrated comparison,
while $\alpha=0.5$ supplies a more concentrated one.
Their realized mean effective numbers are reported rather than assumed to
match the candidate.
The retained populations are transformations of Dirichlet draws rather than
draws from a truncated Dirichlet law.

These laws define comparison populations rather than portfolio strategies that
an investor would implement directly.
Their common feasible set is fully invested, long-only, and subject to a
single-name cap, a recognizable portfolio-constraint template even though the
exact thresholds are stylized rather than rules of a named mandate.
The $12.5\%$ cap rules out portfolios supported on fewer than eight names; the
$15\%$ cap relaxes that minimum to seven.
Changing the cap tests sensitivity to the admissible largest position, while
changing $\alpha$ alters typical concentration within the capped simplex.
Practical mandates may additionally constrain sectors, liquidity, turnover, or
tracking error, none of which these reference laws reproduce.

The named strategy benchmarks answer a different question.
Equal normalized risk weights set $q_i=1/N$ and, through
\eqref{eq:p3-risk-weights}, correspond to inverse-volatility capital weights.
They use marginal volatility but not cross-asset covariance.
The ex post sample GMV instead uses the full return covariance and supplies the
covariance-informed in-sample optimum.
Thus the reference populations locate the news-only allocation within declared
feasible sets, whereas the named benchmarks give its variance an economically
familiar scale.

\subsection{Returns, news, and Wasserstein construction}
\label{sec:p3-data-construction}

The analytical sample follows source, restriction, and final-sample order.
Daily simple returns are Yahoo Finance adjusted-close ratios minus one.
These are raw rather than excess returns; no risk-free series is subtracted.
Exact date alignment supplies \ppnum[0]{n-obs} common observations from
2018-03-19 through 2022-12-30, with missing dates removed rather than imputed.

Nasdaq's per-symbol news archive supplies firm articles from 2018 through 2022.
The analysis applies one ticker-indexed retrieval rule and retains dated article
bodies across the prespecified universe.
This construction is auditable, but the paper neither establishes comprehensive
coverage relative to proprietary news archives nor treats the ticker assignment
as manually validated article-level entity annotation.
The annual coverage screen and balanced-cloud rule then form equal-sized
empirical article distributions for the canonical 2022 geometry.
Eligible articles are ordered deterministically by URL hash before the common
cloud size is imposed, so news volume does not enter the distance merely through
unequal empirical support.
The annual cloud-size ladder appears in \Cref{sec:data-appendix}.

Return alignment leaves \ppnum[0]{n-tickers} firms, and each series is
standardized over the common panel before $\widehat\Sigma_R$ is formed.
This is a coverage-screened survivor panel rather than historical index
membership.

An embedding model maps an article body to a fixed-dimensional numerical vector
and is trained to place semantically related documents nearer in that space.
The primary representation independently encodes each article with the frozen,
4,096-coordinate Qwen3-Embedding-8B bi-encoder
\citep{qwen3_embedding_2025}.
It is obtained through OpenRouter's OpenAI-compatible embeddings API as
\texttt{qwen/qwen3-}\allowbreak\texttt{embedding-8b}; rows are normalized.
This frozen representation is the cross-paper canonical model.
Rather than average a firm's article vectors into one point, the design treats
its row-normalized cloud as an equal-mass empirical distribution.
Balanced quadratic transport then pairs articles across two equal-sized clouds
to minimize their root-mean-square Euclidean chord displacement.
The resulting rooted W$_2$ matrix records the least pairwise semantic
displacement; it does not observe the maintained information-to-exposure
transmission map.
The persisted artifact contains rooted W$_2$ distances, which the allocation
objective above squares once at the certificate boundary.
\Cref{sec:data-appendix} documents the artifact construction and sample roster.

The variance evaluation, $\widehat\Sigma_R$, uses the same 2018--2022 period
as the article geometry after the news-only allocation is constructed.
The exercise is therefore an in-sample descriptive ranking.

\subsection{Allocation-anatomy design}
\label{sec:p3-anatomy-design}

The variance percentile answers whether the news-only allocation lands in an
unusual part of the feasible variance distribution, but not what the rule
selects or which transport separations support that decision.
The allocation-anatomy design therefore reports the canonical firm weights,
aggregates them by sector, and decomposes certificate credit into the additive
unordered terms $q_i q_j W_{2,ij}^{2}$.
These are interpretations of the constructed allocation rather than new
return hypotheses.

\subsection{Expanding-cutoff design}
\label{sec:p3-cutoff-design}

The cutoff diagnostic holds the frozen representation and optimization
rule fixed while expanding the article window through each year-end from 2018
to 2022.
For every cutoff, the analysis reports firm weights, effective $N$, the largest
position, and one-way turnover from the preceding allocation.
The terminal cutoff must reproduce the canonical allocation.
Because the encoder is common across cutoffs and the entire accumulated corpus
changes at each step, this sequence measures decision sensitivity to observed
article coverage; it is neither a point-in-time trading backtest nor an event
study.

\subsection{Representation sensitivity design}
\label{sec:p3-representation-design}

The representation ladder studies sensitivity in the measurement of
information geometry rather than introducing new economic hypotheses or
selecting an encoder from the same return sample.
Every arm holds fixed the priced firms, source article records,
standardized-return covariance, reference laws, caps, and paired bootstrap
schedule while changing one declared representation margin where possible.
Model-specific input limits can still change the effective article text seen by
an encoder.

The first margin is output width.
Qwen3-Embedding-8B is evaluated at its full 4,096 coordinates and at 1,024,
256, and 64 coordinates.
Its Matryoshka training is designed to keep leading-coordinate prefixes
informative, so these arms probe compression within one frozen model rather than
re-estimating coordinates from the return data.
The second margin is model capacity.
The 1,024-coordinate comparison between Qwen3-Embedding-8B and
Qwen3-Embedding-4B holds width fixed, whereas their native 4,096- and
2,560-coordinate comparison changes capacity and width jointly.
Both Qwen models are obtained through OpenRouter.

The third margin is model family.
At 1,024 coordinates, BAAI BGE-large-en-v1.5 changes model family, pretraining,
and effective input length jointly relative to Qwen3-Embedding-8B, so the
comparison is an external sensitivity rather than an identified
architecture effect
\citep{xiao_cpack_2023}.
BGE-large-en-v1.5 is also obtained through OpenRouter.

The final margin is encoder vintage.
The authors' locally trained, 320-coordinate EttaX V0, V1, and V3 encoders hold
architecture, optimization recipe, compute, tokenizer, and token budget fixed
while varying only the Wikipedia training snapshot.
V0 uses 20 December 2017, V1 uses 20 December 2020, and V3 uses 1 August 2026.
Because V3 postdates both the article and return sample, it is a post-sample
negative control rather than a valid point-in-time encoder.
The EttaX encoders also differ materially in capacity from the production Qwen
models, so the vintage contrasts remain descriptive sensitivities rather than
identified vintage effects.
Their benchmark-calibrated equivalence threshold is post-specified, making the
equivalence assessment exploratory.

With the allocation, reference-population estimand, anatomy and cutoff
diagnostics, return evaluation, and four measurement margins defined, the next
section reports the resulting in-sample evidence.

\section{Results}
\label{sec:results}\label{sec:news-variance-rank}

\subsection{Main in-sample variance ranking}

The main result uses the prespecified, 4,096-coordinate
Qwen3-Embedding-8B representation.
Across the four reference laws, its variance percentile ranges from
\ppnum[2]{news-only-percentile-uniform-cap-15-pct}\% to
\ppnum[2]{news-only-percentile-concentrated-cap-15-pct}\%, compared with
\ppnum[1]{equal-percentile-concentrated-cap-15-pct}\% to
\ppnum[1]{equal-percentile-uniform-cap-12p5-pct}\% for equal risk weights.
The variance percentile is the share of feasible reference portfolios with
variance no greater than the candidate's.
Lower is therefore better: a smaller percentile means that fewer reference
portfolios attain equally low or lower variance.
\Cref{tab:news-percentile} reports this ranking under each prespecified
reference law.

\begin{table}[H]
  \centering
  \scriptsize
  \caption{In-sample variance percentile inside four prespecified feasible
    allocation populations. Entries are the share of
    \ppnum[0]{reference-draws} capped draws whose standardized variance is no
  greater than the candidate's. Lower is better.}
  \label{tab:news-percentile}
  \begin{tabular}{lrrr}
    \toprule
    Reference population & Mean effective $N$ & News-only & Equal risk \\
    \midrule
    Uniform, 12.5\% cap
    & \ppnum[2]{reference-effective-n-uniform-cap-12p5}
    & \ppnum[2]{news-only-percentile-uniform-cap-12p5-pct}\%
    & \ppnum[1]{equal-percentile-uniform-cap-12p5-pct}\% \\
    Uniform, 15\% cap
    & \ppnum[2]{reference-effective-n-uniform-cap-15}
    & \ppnum[2]{news-only-percentile-uniform-cap-15-pct}\%
    & \ppnum[1]{equal-percentile-uniform-cap-15-pct}\% \\
    Fixed $\alpha=0.8$, 15\% cap
    & \ppnum[2]{reference-effective-n-effective-n-matched}
    & \ppnum[2]{news-only-percentile-effective-n-matched-pct}\%
    & \ppnum[1]{equal-percentile-effective-n-matched-pct}\% \\
    Concentrated, 15\% cap
    & \ppnum[2]{reference-effective-n-concentrated-cap-15}
    & \ppnum[2]{news-only-percentile-concentrated-cap-15-pct}\%
    & \ppnum[1]{equal-percentile-concentrated-cap-15-pct}\% \\
    \bottomrule
  \end{tabular}
\end{table}

Under the two uniform laws, only
\ppnum[2]{news-only-percentile-uniform-cap-12p5-pct}\% and
\ppnum[2]{news-only-percentile-uniform-cap-15-pct}\% of draws have variance no
greater than the news-only allocation's.
The corresponding shares are
\ppnum[2]{news-only-percentile-effective-n-matched-pct}\% under the fixed
$\alpha=0.8$ law and
\ppnum[2]{news-only-percentile-concentrated-cap-15-pct}\% under the deliberately
more concentrated law.
Thus, a rule constructed without cross-asset return covariance lands near the
first percentile of conventional variance across all four declared feasible
populations.

The two concentrated populations make this ranking more informative than the
uniform comparisons alone, but they do not furnish an effective-$N$ match.
The fixed $\alpha=0.8$ law has mean effective $N$ of
\ppnum[2]{reference-effective-n-effective-n-matched}, compared with the
candidate's \ppnum[2]{news-only-effective-number-assets}, and only
\ppnum[2]{news-only-percentile-effective-n-matched-pct}\% of its draws attain
equally low or lower variance.
Under the more concentrated law, mean effective $N$ falls to
\ppnum[2]{reference-effective-n-concentrated-cap-15}, and the share is only
\ppnum[2]{news-only-percentile-concentrated-cap-15-pct}\%.
The low rank therefore survives comparisons that materially alter the
concentration profile within the declared feasible class.
Because both reference laws remain more diffuse than the candidate, this
sensitivity does not eliminate concentration as an explanation of the rank.

Equal risk weights provide a familiar point of comparison within the same laws.
Their percentiles are
\ppnum[1]{equal-percentile-uniform-cap-12p5-pct}\%,
\ppnum[1]{equal-percentile-uniform-cap-15-pct}\%,
\ppnum[1]{equal-percentile-effective-n-matched-pct}\%, and
\ppnum[1]{equal-percentile-concentrated-cap-15-pct}\%, respectively.
These ranks remain well above the news-only ranks under every reference law.
The comparison is a descriptive location within prespecified portfolio
populations, not a hypothesis test or evidence of prospective performance.

\subsection{Portfolio characteristics and conventional benchmarks}
\label{sec:news-variance-benchmark}

The percentile result says where the candidate lies within feasible reference
populations but not how far it sits from familiar portfolio strategies.
The equal-risk benchmark sets $q_i=1/N$.
Using the mapping $q_i(x)=x_i\sigma_i/A(x)$ gives
\[
  q_i=\frac1N
  \quad\Longleftrightarrow\quad
  x_i^{\mathrm{IV}}
  =\frac{\sigma_i^{-1}}{\sum_j\sigma_j^{-1}}.
\]
It is therefore the inverse-volatility capital portfolio expressed in the
normalized risk-weight coordinates of the empirical exercise, not an arbitrary
second equal-weight portfolio.
The sample GMV provides the opposite benchmark: it uses the complete in-sample
covariance matrix that the news-only construction excludes.
\Cref{tab:news-variance-benchmark} provides that calibration using the same
standardized return panel.

\begin{table}[H]
  \centering
  \scriptsize
  \caption{In-sample descriptive portfolio benchmarks. Standardized variance
    is computed on the complete 2018--2022 aligned return panel; relative
    variance indexes the long-only sample GMV to 100. Returns evaluate the
  news-only allocation after construction but do not determine its weights.}
  \label{tab:news-variance-benchmark}
  \begin{tabular}{lrrrr}
    \toprule
    Portfolio & Standardized variance & Rel. to GMV & Max. weight
    & Effective $N$ \\
    \midrule
    News-only
    & \ppnum[3]{news-only-variance}
    & \ppnum[1]{news-only-gmv-relative-variance-pct}
    & \ppnum[1]{news-only-maximum-weight-pct}\%
    & \ppnum[2]{news-only-effective-number-assets} \\
    Equal risk weights
    & \ppnum[3]{equal-standardized-variance}
    & \ppnum[1]{equal-gmv-relative-variance-pct}
    & -- & -- \\
    Sample GMV
    & \ppnum[3]{sample-gmv-variance}
    & 100.0 & -- & -- \\
    \bottomrule
  \end{tabular}
\end{table}

The news-only standardized variance is
\ppnum[1]{news-only-equal-variance-reduction-pct}\% lower than the variance of
the equal-risk benchmark.
It nevertheless remains
\ppnum[1]{news-only-gmv-excess-variance-pct}\% above the ex post long-only
sample GMV.
The empirical claim is therefore not that information geometry solves the
conventional GMV problem without a covariance matrix.
Rather, the news-only allocation lands unusually low in the variance
distribution generated by the declared feasible laws while remaining short of
the covariance-informed in-sample optimum.

The concentration diagnostics describe the decision object behind that rank.
Its largest holding is \ppnum[1]{news-only-maximum-weight-pct}\%, and its
effective number of names is
\ppnum[2]{news-only-effective-number-assets}.
The candidate therefore differs from equal risk weighting and is not reduced to
a one-name solution.
Together with the concentration-adjusted reference comparisons, these
diagnostics show that the low rank survives laws that place more mass on
concentrated portfolios.
Because none matches the candidate's realized effective number of names, they
do not eliminate scalar concentration as a confound.
They do not establish stability of the individual holdings or performance
outside this sample.

\subsection{Anatomy of the news-only allocation}
\label{sec:p3-portfolio-anatomy}

The scalar concentration diagnostics do not reveal which firms the rule
selects or which parts of the information geometry create its certificate.
\Cref{fig:p3-portfolio-anatomy} opens that decision object at three levels:
firm weights, additive firm-pair credits, and sector weights.

The pair layer is the informative attribution level for this optimum.
To see why, let $M=(W_{2,ij}^{2})_{i,j}$ and write the zero-slack objective as
$\mathcal C(q)=\tfrac12q^{\top}Mq$.
For every firm in the positive support of this solution, whose unit cap is
nonbinding, the simplex first-order condition gives
$(Mq^{\star})_i=\lambda$.
Multiplying by $q_i^{\star}$ and summing over the support yields
$\lambda=q^{\star\top}Mq^{\star}$.
Consequently, the natural firm-level share proposed by the quadratic
decomposition satisfies
\begin{equation}\label{eq:p3-firm-credit-identity}
  \frac{q_i^{\star}(Mq^{\star})_i}
  {q^{\star\top}Mq^{\star}}=q_i^{\star}
  \qquad\text{whenever }q_i^{\star}>0.
\end{equation}
Firm certificate shares therefore reproduce the portfolio weights rather than
provide a second diagnostic.
The unordered pair terms remain informative:
\begin{equation}\label{eq:p3-pair-credit-decomposition}
  c_{ij}:=q_i^{\star}q_j^{\star}W_{2,ij}^{2},
  \qquad
  \mathcal C(q^{\star})=\sum_{i<j}c_{ij}.
\end{equation}
They identify which weighted separations the optimizer actually uses.

\begin{figure}[H]
  \centering
  \ppfigure{0.95}{portfolio_anatomy}
  \caption{Anatomy of the canonical news-only allocation. Panel (a)
    orders firms by sector and then by weight; horizontal lines mark the equal
    risk weight and the \ppnum[1]{reference-tight-cap-pct}\% cap used by the
    tightest reference population.
    Alternating ticker positions are labelled for legibility; all
    \ppnum[0]{n-tickers} firm weights remain plotted.
    Panel (b) reports the twelve largest unordered terms $c_{ij}$ in
    \eqref{eq:p3-pair-credit-decomposition} as shares of total certificate
    credit and distinguishes within- from cross-sector pairs. Panel (c)
    compares sector weights with the sector shares induced by equal risk
    weights. All panels describe the full-sample news-only allocation; they do
  not use returns or assign a causal role to sector membership.}
  \label{fig:p3-portfolio-anatomy}
\end{figure}

The solution assigns positive weight to
\ppnum[0]{anatomy-active-firms} of the \ppnum[0]{n-tickers} firms, so the
effective-$N$ statistic reflects a portfolio with many small positions and a
limited set of larger ones rather than diffuse weight on every name.
\Cref{tab:p3-top-holdings} makes the largest positions tangible.
For each row, the strongest partner is the holding that produces its largest
pair term in \eqref{eq:p3-pair-credit-decomposition}.
A large position need not have the largest pair credit: the latter also
depends on the counterpart's weight and their squared transport separation.

% AUTO-GENERATED by pipeline render-paper3-portfolio-anatomy
% -- do not edit by hand.
\begin{table}[H]
\centering
\scriptsize
\setlength{\tabcolsep}{3.0pt}
\begin{tabularx}{\linewidth}{lXlrrr}
\toprule
Ticker & Firm & Sector & Weight (\%) & Partner & Pair credit (\%) \\
\midrule
AZN & AstraZeneca PLC & Healthcare & 10.86 & GOOG & 1.34 \\
GOOG & Alphabet Inc. & Communication Services & 5.31 & AZN & 1.34 \\
NFLX & Netflix Inc. & Communication Services & 4.83 & AZN & 1.21 \\
XEL & Xcel Energy Inc. & Utilities & 4.72 & AZN & 1.20 \\
TSCO & Tractor Supply Company & Consumer Cyclical & 4.21 & AZN & 0.97 \\
UAL & United Airlines Holdings Inc. & Industrials & 4.18 & AZN & 0.99 \\
BMRN & BioMarin Pharmaceutical Inc. & Healthcare & 4.02 & AZN & 0.92 \\
TMUS & T-Mobile US Inc. & Communication Services & 3.46 & AZN & 0.86 \\
\bottomrule
\end{tabularx}
\caption{Largest positions in the canonical news-only allocation. Partner identifies the holding that forms the largest unordered additive certificate term with the row firm; pair credit reports that term as a percentage of total certificate credit.}
\label{tab:p3-top-holdings}
\end{table}


The sector panels supply a diagnostic rather than an explanation.
The news-only sector weights differ visibly from the sector shares implied by
equal risk weights.
Across the additive pair terms,
\ppnum[1]{anatomy-cross-sector-credit-pct}\% of certificate credit is
cross-sector and \ppnum[1]{anatomy-within-sector-credit-pct}\% is
within-sector.
Thus the reported certificate is not generated by within-sector separation
alone.
The split does not establish that transport geometry contains information
beyond sector labels: sector sizes, portfolio weights, and the number of
available cross-sector pairs also affect it.
It locates the certificate credit while leaving the economic source of those
separations open.

\subsection{Evolution across expanding information cutoffs}
\label{sec:p3-portfolio-evolution}

The canonical allocation pools articles through 2022.
To examine whether that decision is peculiar to the terminal corpus, the same
zero-slack problem is re-solved after expanding the article window through
each year-end from 2018 to 2022.
\Cref{fig:p3-portfolio-evolution} shows the resulting sequence.
These are expanding information cutoffs, not separate calendar-year article
samples.

\begin{figure}[H]
  \centering
  \ppfigure{0.95}{portfolio_evolution}
  \caption{Evolution of the news-only allocation under expanding
    article cutoffs. Panel (a) reports firm weights for article windows ending
    in 2018, 2019, 2020, 2021, and 2022, with firms grouped by sector.
    Alternating ticker positions are labelled; every firm remains in each
    heatmap column.
    Panel (b) reports the effective number of names, the largest weight and its
    ticker, and one-way turnover
    $\tfrac12\lVert q_t-q_{t-1}\rVert_1$. The common frozen encoder is held
    fixed, and the terminal 2022 column reproduces the canonical allocation.
    The sequence measures decision sensitivity to accumulated articles; it is
  neither a prospective return test nor an event-study design.}
  \label{fig:p3-portfolio-evolution}
\end{figure}

The portfolio changes as the information set grows without becoming either
one-name concentrated or equal-risk weighted.
Across the five cutoffs, effective $N$ remains between
\ppnum[2]{anatomy-vintage-effective-n-min} and
\ppnum[2]{anatomy-vintage-effective-n-max}, while the largest weight ranges
from \ppnum[1]{anatomy-vintage-maximum-weight-min-pct}\% to
\ppnum[1]{anatomy-vintage-maximum-weight-max-pct}\%.
One-way turnover ranges from \ppnum[1]{anatomy-turnover-min-pct}\% to
\ppnum[1]{anatomy-turnover-max-pct}\%.
The largest reallocation occurs when the first one-year corpus is extended
through 2019:
\ppnum[1]{anatomy-turnover-2018-2019-pct}\%, compared with
\ppnum[1]{anatomy-turnover-2019-2020-pct}\% from 2019 to 2020.
The sequence therefore does not display a uniquely large break at the 2020
cutoff.
It cannot isolate a COVID effect in any case, because the design changes the
entire accumulated article distribution at once and holds a common encoder
fixed across cutoffs.

\subsection{Representation sensitivity}
\label{sec:p3-representation-results}

The expanding-cutoff sequence changes the information set while holding the
representation fixed.
The final empirical question changes the representation itself.
\Cref{tab:p3-representation-main} shows the four measurement margins;
\Cref{sec:p3-representation-ablation} reports the full ladder and paired
bootstrap contrasts.

% AUTO-GENERATED -- do not edit by hand.
\begin{table}[H]
\centering
\scriptsize
\caption{Selected representation sensitivity of the news-only allocation. The percentile column reports the percentage of allocations under the fixed $\alpha=0.8$ reference law whose standardized variance is no greater than the candidate's. Relative variance indexes the long-only sample GMV to 100. Lower is better in both columns. The complete ladder and paired bootstrap contrasts appear in \Cref{sec:p3-representation-ablation}.}
\label{tab:p3-representation-main}
\setlength{\tabcolsep}{7.0pt}
\begin{tabular}{lrrr}
\toprule
Encoder & Dim. & Fixed-$\alpha=0.8$ percentile & Rel. GMV \\
\midrule
Qwen3-8B full (primary) & 4096 & 0.19 & 147.5 \\
Qwen3-8B@1024 & 1024 & 0.08 & 146.7 \\
Qwen3-4B@1024 & 1024 & 0.29 & 148.5 \\
BGE-large-v1.5 full & 1024 & 0.03 & 145.3 \\
Qwen3-8B@64 & 64 & 0.01 & 144.6 \\
\addlinespace[1.5pt]
EttaX V0 & 320 & 0.53 & 149.7 \\
EttaX V1 & 320 & 1.59 & 152.0 \\
EttaX V3 & 320 & 1.20 & 151.4 \\
\bottomrule
\end{tabular}
\end{table}


Moderate within-model compression leaves the outcome nearly unchanged:
Qwen3-8B at 1,024 coordinates is close to the full-width primary representation
in both reported metrics, and their paired interval includes zero.
At the same fixed width, Qwen3-4B ranks higher in conventional variance than
Qwen3-8B, but their paired difference does not survive Holm adjustment.
BGE-large at 1,024 coordinates is again close to Qwen3-8B, so the sample does
not support a simple model-family ordering.

The width ladder is non-monotone.
The 64-coordinate Qwen3-8B stress cell has the lowest conventional-variance
rank in the ladder, but its contrast with the 256-coordinate cell does not
survive Holm adjustment.
That outcome does not make 64 coordinates the preferred specification: the
full-width 8B representation was prespecified, and every arm is evaluated on
the same return sample.
Instead, its point ranking shows that extreme compression can materially alter
the geometry and selected allocation, in this case in a direction that
improves the in-sample variance benchmark without an adjusted rejection of
zero difference against the 256-coordinate cell.

The matched EttaX vintages supply a separate sensitivity.
Their point outcomes vary across the three training snapshots, but none of the
three contrasts survives Holm adjustment.
The post-specified equivalence criterion is met only for V3--V1; the intervals
comparing V0 with V1 or V3 are too wide to establish equivalence.
Because V3 postdates the sample and the equivalence threshold is exploratory,
these comparisons neither identify a causal vintage effect nor establish broad
representation invariance.
Taken together, the ladder does not support a monotone ``larger model is
better'' account.
It shows instead that the primary result is empirically strong within the
declared benchmark but remains sensitive to the measurement layer that creates
the information geometry.

These comparisons are descriptive and in-sample: the same 2018--2022 return
panel evaluates the news-only allocations after their construction.
They describe how the rule ranks against the declared reference populations,
not how it will perform prospectively.
\Cref{sec:limitations} interprets what these results establish within the
maintained framework and where their boundaries lie.
% Temporarily omitted while the chronological-validation
% implementation is audited.
\iffalse
\subsection{Chronological certificate validation}
\label{sec:chronological-validation}

No prespecified transmission cell clears the calibration coverage rule.
Following the frozen protocol, the evaluation therefore uses zero credit and
the benchmark of perfect dependence rather than selecting a more favourable
holdout cell.  This is a validation failure for an informative certificate,
not a reason to change the horizon, grid, or transmission specification.
\Cref{sec:calibration-boundary} states precisely which region that failure
covers, because a rule the declared grid never reaches is a weaker finding than
a region the data reject.

\Cref{tab:validation-headline} reports the evidence in its decision order:
validity first, then informativeness, then portfolio value.  Coverage is shown
with its one-sided 95\% lower endpoint in brackets.  Daily future variances are
reported on the raw return scale.  Required $\delta$ is the mean additional
normalized residual-covariance budget needed by the operational inverse-
volatility bound; it is a sensitivity diagnostic, not an estimated structural
parameter.

\begin{table}[H]
  \centering
  \caption{Frozen chronological certificate validation.  Coverage entries are
    rates with one-sided 95\% lower endpoints in brackets.  Credit and floor
  activation are percentages; future variance is daily raw return variance.}
  \label{tab:validation-headline}
  \begin{tabular}{llr}
    \toprule
    Order & Outcome & Estimate \\
    \midrule
    Validity & Operational coverage
    & \ppnum[1]{validation-operational-coverage-pct}\%
    [\ppnum[1]{validation-operational-coverage-lower-pct}\%] \\
    Validity & Geometry-only coverage
    & \ppnum[1]{validation-geometry-only-coverage-pct}\%
    [\ppnum[1]{validation-geometry-only-coverage-lower-pct}\%] \\
    Informativeness & Mean lower certificate credit
    & \ppnum[1]{validation-mean-lower-credit-pct}\% \\
    Informativeness & Active pairwise floors
    & \ppnum[1]{validation-floor-activation-pct}\% \\
    Decision & Certificate-portfolio future variance
    & \ppnume[2]{validation-variance-certificate} \\
    Decision & Inverse-volatility future variance
    & \ppnume[2]{validation-variance-inverse-volatility} \\
    Decision & Certificate minus inverse volatility
    & \ppnume[2]{validation-variance-difference} \\
    Sensitivity & Mean required $\delta$
    & \ppnum[6]{validation-required-delta} \\
    \bottomrule
  \end{tabular}
\end{table}

Because calibration activates the fallback, point and lower-endpoint credit are
both zero in the holdout.  Operational and geometry-only coverage consequently
test the perfect-dependence benchmark rather than an informative deduction.
The certificate portfolio also coincides with inverse volatility under the
fallback, so their future-variance difference is mechanically zero.  The
sample-GMV and Ledoit--Wolf rows remain governed secondary comparisons, but
cannot rescue decision value for a certificate that calibration did not admit.

\Cref{fig:certificate-validation} plots certified diversification against
subsequently realized diversification for 500 fixed-seed capped random
portfolio identities, averaging each identity across the 40 holdout formations.
The horizontal axis is the certified deduction from perfect dependence; the
vertical axis is the realized deduction computed with the same formation-time
volatility budget.  Points on or above the dashed 45-degree line satisfy the
bound.  Under the frozen fallback all points have zero certified deduction,
making the figure a transparent display of the failed informativeness gate
rather than an ex post refit.

\begin{figure}
  \centering
  \ppfigure{0.72}{certificate_validation}
  \caption{Certified versus subsequently realized diversification.  Each point
    is the chronological mean for one of 500 seed-42 random portfolio
    identities across 40 monthly holdout formations.  The dashed 45-degree line
    is the bound benchmark.  The vertical pile at zero reflects the
    preregistered zero-credit fallback after no calibration cell passed; it is
  not evidence that text geometry explains realized risk.}
  \label{fig:certificate-validation}
\end{figure}

\subsection{What the declared grid could test}\label{sec:calibration-boundary}

The two calibration conditions do not fail together.  Of the
\ppvalue{calibration-cells-total} cells,
\ppvalue{calibration-cells-margin-pass} satisfy the mean-margin condition and
\ppvalue{calibration-cells-passing} satisfy both.  Coverage alone binds: the
bound is conservative on average across most of the grid and still misses the
required share of calibration formations.

\Cref{fig:calibration-coverage} shows where that leaves the design.  Coverage
increases monotonically in both $L$ and $\tau$, so the best-covering cell is the
most conservative one the grid contains, at
$L=\ppnum[1]{calibration-best-coverage-scale}$ and
$\tau=\ppnum[2]{calibration-best-coverage-slack}$ --- exactly the corner of the
declared $L\times\tau$ range.  It reaches
\ppnum[1]{calibration-best-coverage-pct}\% coverage with a
\ppnum[1]{calibration-best-coverage-lower-pct}\% lower endpoint against a 90\%
rule, while still carrying
\ppnum[1]{calibration-best-coverage-credit-pct}\% mean lower credit.

The correct reading is therefore narrower than it may first appear.  The
evidence establishes that the prespecified transmission region does not reach
the declared coverage standard, not that no transmission region could.  Because
coverage is still rising at the grid boundary, a more conservative region
outside the declared range may satisfy the rule; the credit it could carry would
be correspondingly smaller, since credit falls as $L$ and $\tau$ rise
(\Cref{fig:certificate-surface}).  We do not estimate where that crossing lies.
Extending the grid after observing this outcome would convert a preregistered
design into a tuned one, which is the specific failure the protocol exists to
prevent.  Locating the crossing requires a fresh preregistration, and the
paper's conclusions are stated for the region actually declared.

\begin{figure}[H]
  \centering
  \ppfigure{0.72}{calibration_coverage}
  \caption{Calibration coverage lower endpoint across the declared
    transmission grid.  Each line fixes $\tau$ and varies $L$; the dashed line
    is the 90\% calibration rule.  Coverage rises with conservatism and is
    highest at the grid corner, so the rule is never reached inside the declared
  range rather than being approached and rejected.}
  \label{fig:calibration-coverage}
\end{figure}

Ticker, within-sector, stale-2018-geometry, and 1,000 label-permutation placebos
are retained in the governed results.  Under the zero-credit fallback they are
necessarily uninformative: relabelling or ageing a geometry that receives zero
credit cannot change the bound.  Turnover, concentration, and article-draw
weight stability are likewise diagnostics rather than headline outcomes.  The
residual sensitivity budgets $\delta\in\{0,0.025,0.05\}$ only show how an
assumed additive residual allowance changes coverage; they do not reinterpret
credit against perfect dependence as explained realized risk.
\fi
